\begin{table}[H]
  \caption{Residue landmarks and numbering used across the six-enzyme panel.}
  \label{tab:residue_correspondence}
  \scriptsize
  \setlength{\tabcolsep}{4pt}
  \begin{tabularx}{\linewidth}{@{}ll>{\raggedright\arraybackslash}Xll>{\raggedright\arraybackslash}X@{}}
    \toprule
    Enzyme & Catalytic triad & Oxyanion, subsite, or interface landmarks & W-loop Trp & Adjacent pair & Numbering note \\
    \midrule
    IsPETase & S160/D206/H237 & Y87 and M161 backbone-N donors & W185 & S214/I218 & Simulation numbering. Trp156 in 5XG0 maps to W185 in the full-length simulation model. \\
    TfCut1 & S131/D177/H209 & Y61 and M132 backbone-N donors & W156 & H185/F189 & Simulation numbering. \\
    LCC & S165/D210/H242 & Y95 and the M166 region & W190 & H218/F222 & Simulation numbering. \\
    FoCut5a & S122/D177/H190 & Flap 75--94 and binding-loop region 173--193 & --- & --- & Simulation numbering is $+1$ relative to 5AJH and the literature numbering in these mapped regions. \\
    HiC & S105/D160/H173 & S28/T29 and F70 cleft sites; L167/I168/I169 catalytic-His-side rim & --- & --- & Simulation numbering. \\
    PHL7 & S131/D177/H209 & F63/M132/I179 subsite-I landmarks & W156 & H185/F189 & Simulation and PHL7 reference numbering coincide at the listed positions. \\
    \bottomrule
  \end{tabularx}
  \par\vspace{2pt}
  \begin{minipage}{\linewidth}
    \footnotesize
    Residue identifiers in the manuscript use the local simulation numbering. Dashes indicate that no W-loop tryptophan or adjacent-pair landmark was assigned in this study.
  \end{minipage}
\end{table}
